% This template was partially made with assistance from Microsoft Co-pilot

% ------------------------------------------------------------------
% CAPSTONE SOFTWARE REQUIREMENTS SPECIFICATION (SRS)

% 
% This template is inspired by IEEE 29148 but simplified for use in 
% a senior software engineering capstone project. 
% 
% The goal is not to maximize document length. 
% The goal is to produce a clear, testable, and traceable set of 
% requirements that can guide design, implementation, and testing. 
% 
% Recommended workflow: 
% Stakeholder Requirements 
% -> 
% Use Cases 
% -> 
% Functional Requirements 
% -> 
% Acceptance Tests (will be in your VnV plan)
% -> 
% Traceability Matrix (will excluded traceability for tests, that will be in the VnV plan)
% 
% Students should update this document throughout the project. 
% ------------------------------------------------------------------

\documentclass[12pt]{article}

\usepackage{amsmath, mathtools}
\usepackage{amsfonts}
\usepackage{amssymb}
\usepackage{graphicx}
\usepackage{colortbl}
\usepackage{xr}
\usepackage{hyperref}
\usepackage{longtable}
\usepackage{xfrac}
\usepackage{tabularx}
\usepackage{float}
\usepackage{siunitx}
\usepackage{booktabs}
\usepackage{caption}
\usepackage{pdflscape}
\usepackage{afterpage}
\usepackage{parskip}

\hypersetup{
    colorlinks=true,       % false: boxed links; true: colored links
    linkcolor=red,          % color of internal links (change box color with linkbordercolor)
    citecolor=green,        % color of links to bibliography
    filecolor=magenta,      % color of file links
    urlcolor=cyan           % color of external links
}

\setlength{\parindent}{0pt}

\input{../Comments.text}
\input{../Common.text}

% Used so that cross-references have a meaningful prefix

\newcounter{connum} %System Constraint (C) Number
\newcommand{\conref}[1]{C.\ref{#1}}

\newcounter{srnum} %Stakeholder Requirement (SR) Number
\newcommand{\srthesrnum}{SR\thesrnum}
\newcommand{\srref}[1]{SR\ref{#1}}

\newcounter{ucnum} %Use Case (UC) Number
\newcommand{\uctheucnum}{UC\theucnum}
\newcommand{\ucref}[1]{UC\ref{#1}}

\newcounter{reqnum} %Functional Requirement (FR) Number
\newcommand{\rthereqnum}{FR\thereqnum}
\newcommand{\rref}[1]{FR\ref{#1}}

\newcounter{nfrnum} %Non-Functional Requirement (NFR) Number
\newcommand{\rthenfrnum}{NFR\thenfrnum}
\newcommand{\nfrref}[1]{NFR\ref{#1}}

\newcounter{lcnum} %Likely Change (LC) Number
\newcommand{\lthelcnum}{LC\thelcnum}
\newcommand{\lcref}[1]{LC\ref{#1}}

\usepackage{fullpage}

\begin{document}

\title{Software Requirements Specification for \progname: An Ophthalmology Preparation Game for Children} 
\author{\authname}
\date{\today}
	
\maketitle

~\newpage

\pagenumbering{roman}

\tableofcontents

~\newpage

\begin{table}[hp]
\caption{Revision History} \label{TblRevisionHistory}
\begin{tabularx}{\textwidth}{llX}
\toprule
\textbf{Date} & \textbf{Developer(s)} & \textbf{Change}\\
\midrule
2026/10/12 & Mohammad Bilal & Completed sections 1.1, 1.2, 2.1, \& contributed to 1.4 \\
2026/10/13 & Ahsan Muzammil & Completed sections 2.2, 5, 7, \& contributed to Appendix D, E \\
\bottomrule
\end{tabularx}
\end{table}

\newpage

\pagenumbering{arabic}

% ============================================================== 
\section{Introduction} 
% ==============================================================

This document is the Software Requirements Specification (SRS) for \progname{}, a mobile game
that prepares children aged 4 to 7 years for paediatric ophthalmology appointments.
This document defines the expected behaviour, scope, and requirements of \progname{},
establishing a basis for implementation, verification, and validation by describing what the system
is expected to accomplish and the conditions under which it will be used.

This document follows the IEEE SRS template provided in the course repository \cite{SmithCapTemplate},
with guidance from ISO/IEC/IEEE 29148 \cite{ISOIECIEEE291482018} for requirements engineering and
ISO/IEC/IEEE 15288 \cite{ISOIECIEEE152882023} for systems life-cycle management. The team made
modifications to the original template to better suit the project's needs. In Section~\ref{non-func-req}
(Non-Functional Requirements), the requirements are organized according to the Volere template provided
in the course repository \cite{SmithCapTemplate}. The AI and Machine Learning Requirements section was
also removed, as it is not applicable to the project.

\subsection{Purpose}

The purpose of \progname{} is to reduce anxiety and uncertainty in children before paediatric
ophthalmology appointments by providing accessible, age-appropriate preparation for common eye
examination procedures. \progname{} aims to improve children's preparedness and willingness to participate while helping
parents and healthcare professionals prepare them for appointments. Intended outcomes include
reduced anxiety, improved cooperation, sustained engagement across the target age range,
more effective preparation at home, and greater clinical efficiency. The application complements,
rather than replaces, professional guidance and medical examinations.

\subsection{Scope} \label{sec_Scope}

\progname{} is a mobile application intended to help children aged 4 to 7 years become familiar
with common paediatric ophthalmology procedures before attending an appointment. The application
will present preparation activities as age-appropriate interactive games that allow children to
practice actions and learn what they may experience during an eye examination.

The primary users are children preparing for ophthalmology appointments, with parents or guardians
supporting their use of the application. Ophthalmologists and child-life specialists are secondary
stakeholders who may provide clinical and developmental guidance during development and may use
the application as part of a broader preparation process.

The application will provide interactive activities corresponding to common ophthalmology
procedures, provide feedback during activities, track completion of preparation activities, and
allow users to repeat activities when needed. The application will be designed to support
preparation before an appointment, including preparation in a home environment.

The expected benefits are increased familiarity with ophthalmology procedures, reduced uncertainty
and anxiety, improved cooperation during examinations, sustained engagement, and potentially
shorter or less difficult clinical appointments.

\progname{} will not diagnose medical conditions, provide medical advice, replace an ophthalmologist,
or determine whether a child requires a particular medical procedure. The application will not
replace an actual clinical examination. Medical decisions will remain the responsibility of qualified
healthcare professionals.

The initial system will focus on preparation for common ophthalmology procedures and will not
attempt to cover every possible procedure or clinical situation. Features identified as stretch
goals, such as augmented reality simulations and personalized preparation based on appointment
information, are outside the minimum required scope of the initial system.

\subsection{Stakeholders}
\begin{itemize}
\item Customer
\item End Users
\item Project Sponsor
\item Development Team
\item Administrators
\end{itemize}

\subsection{Definitions and Acronyms} \label{sec_DefinitionsAndAcronyms}

\renewcommand{\arraystretch}{1.3}

\begin{longtable}{@{}p{0.35\textwidth} p{0.62\textwidth}@{}}
\caption{Definitions and Acronyms}
\label{tab:DefinitionsAndAcronyms}\\

\toprule
\textbf{Term} & \textbf{Definition} \\
\midrule
\endfirsthead

\multicolumn{2}{c}{\tablename\ \thetable{} -- continued from previous page} \\
\toprule
\textbf{Term} & \textbf{Definition} \\
\midrule
\endhead

\midrule
\multicolumn{2}{r}{Continued on next page} \\
\endfoot

\bottomrule
\endlastfoot

Software Requirements Specification (SRS) &
A document defining a software system's requirements, scope, behaviour, and constraints. \\

\progname{} &
The application described in this document for preparing children for paediatric ophthalmology appointments. \\

Paediatric Ophthalmology &
The medical specialty focused on children's eye health and conditions. \\

Ophthalmologist &
A physician specializing in eye health and eye conditions. \\

Child-Life Specialist &
A healthcare professional who helps children cope with medical experiences and procedures. \\

Eye Examination &
A clinical assessment of eye health and vision. \\

Eye Drops &
Liquid medication applied to the eyes, including to dilate pupils for examination. \\

Slit-Lamp Examination &
An eye examination using a microscope and a narrow beam of light. \\

Cycloplegic Eye Examination &
An eye examination using drops to relax focusing muscles and assess refractive error. \\

Functional Requirement (FR) &
A specification of a capability or behaviour the application must provide. \\

Non-Functional Requirement (NFR) &
A specification of the application's quality attributes or constraints. \\

Stakeholder Requirement (SR) &
A statement of a stakeholder's needs or expectations. \\

Use Case (UC) &
A description of how an actor uses the application to achieve a goal. \\

Verification and Validation (V\&V) &
Processes for checking requirement compliance and suitability for user needs. \\

Requirements Traceability Matrix (RTM) &
A table linking stakeholder requirements, use cases, and functional and non-functional requirements. \\

Augmented Reality (AR) &
Technology that overlays digital content onto the real world. \\

\end{longtable}

\subsection{Characteristics of Intended Reader} \label{sec_IntendedReader}

The intended readers are the individuals who will read, review and maintain this document. Their characteristics are:
\begin{itemize}
\item An understanding of requirements engineering, at least at the level of an upper-year undergraduate course in a software
engineering/computer science program. This is so the reader can understand the meaning behind the structure (e.g., Personas, Scenarios, etc.)
and terminology (e.g., Functional vs. Non-Functional Requirements) used.
\item An understanding of discrete mathematics, at least at the level of an undergraduate course in a software engineering/computer science program.
This is so the reader can understand the notation used in formal specifications of the system's behaviour.
\item Familiarity with touchscreens and mobile applications/games, at a general level expected of someone whose interacted
with such applications as a consumer.
\item Familiarity with the problem the system is attempting to address, at a level which can be achieved from reading
the Problem Statement document.
\end{itemize}

Knowledge of clinical terms and processes are not a pre-requisite for the intended reader. Such terminology and processes are defined in
Section~\ref{sec_DefinitionsAndAcronyms} and Section~\ref{sec_Scope}, respectively.

% ============================================================== 
\section{System Overview}

This section provides general information about the system.  It identifies the
problem, the vision, and the interfaces between the system and its environment.
This section also describes user characteristics and lists the system
constraints.

\subsection{Problem Statement}

Paediatric ophthalmology appointments can be stressful and unfamiliar for children.
Procedures such as vision testing, eye drops, and slit-lamp examinations may cause anxiety, particularly when children do not know what to expect.
This anxiety can make it difficult to cooperate during examinations, potentially affecting the quality and completeness of clinical assessments.
For example, a study of 298 children undergoing cycloplegic eye examinations found that 26\% experienced significant distress before receiving eye drops,
while uncooperative children were also more anxious and required more time to complete the examination \cite{ChiaEtAl2015}.

Children and their families therefore need adequate preparation for common ophthalmology procedures before attending an appointment.
Without appropriate preparation, unfamiliar clinical experiences may contribute to increased anxiety, reduced cooperation, and longer or more difficult appointments.
This project addresses the need for an accessible, child-friendly way to help children become more familiar and comfortable with the paediatric ophthalmology experience.

\subsection{Product Vision}

\progname{} is an interactive mobile application designed to help
children aged 3--7 feel more comfortable and prepared for paediatric
ophthalmology appointments. Through age appropriate games and activities,
the application introduces children to common procedures, including
vision testing, eye drop administration, and slit lamp examinations.
By allowing children and their families to become familiar with these
procedures before their appointments, \progname{} aims to reduce
anxiety, improve cooperation during examinations, and support a
more efficient clinical experience.

The primary goals of \progname{} are:

\begin{itemize}
    \item \textbf{Reduce Children's Anxiety:}
    Help children understand what to expect during common ophthalmology
    procedures, reducing anxiety associated with unfamiliar clinical
    experiences.

    \item \textbf{Improve Children's Cooperation During Examinations:}
    Prepare children to participate more cooperatively in clinical
    examinations, reducing interruptions and repeated attempts
    during procedures.

    \item \textbf{Maintain Engagement Across Age Groups:}
    Provide engaging and age appropriate activities that encourage
    children across the target age range to actively participate
    in their preparation.

    \item \textbf{Enable Effective At-Home Preparation:}
    Allow children and their families to independently prepare
    for upcoming ophthalmology appointments through accessible
    and understandable activities.

    \item \textbf{Improve Clinical Efficiency:}
    Support more efficient ophthalmology examinations by improving
    children's preparedness and reducing difficulties associated
    with completing clinical procedures.
\end{itemize}

These goals establish the intended outcomes of \progname{} and provide
the foundation for the functional and non-functional requirements
defined in this document.
  
\subsection{System Context}

\plt{Include users, external systems, APIs, databases, and third-party services.}

\plt{Your system context will include a figure that shows the abstract view of
  the software.  In some cases the diagram will show other external
  entities, besides the user.  For instance, when the software product is a
  library, the user will be another software program, not an actual end user.
  If there are system constraints that the software must work with external
  libraries, these libraries can also be shown on the System Context diagram.
  They should only be named with a specific library name if this is required by
  the system constraint.}

\begin{figure}[h!]
\begin{center}
 \includegraphics[width=0.6\textwidth]{SystemContextFigure}
\caption{System Context}
\label{Fig_SystemContext} 
\end{center}
\end{figure}

\plt{For each of the entities in the system context diagram its responsibilities
  should be listed.  Whenever possible the system should check for data quality,
  but for some cases the user will need to assume that responsibility.  The list
  of responsibilities should be about the inputs and outputs only, and they
  should be abstract.  Details should not be presented here.  However, the
  information should not be so abstract as to just say ``inputs'' and
  ``outputs''.  A summarizing phrase can be used to characterize the inputs.
  For instance, saying ``material properties'' provides some information, but it
  stays away from the detail of listing every required properties.}

\begin{itemize}
\item User Responsibilities:
\begin{itemize}
\item 
\end{itemize}
\item \progname{} Responsibilities:
\begin{itemize}
\item Detect data type mismatch, such as a string of characters instead of a
  floating point number
\item 
\end{itemize}
\end{itemize}

\plt{Identify in what context the software will typically be used.  Is it for
exploration? education? engineering work? scientific work?. Identify whether it
will be used for mission-critical or safety-critical applications.} 

\plt{This additional context information is needed to determine how much effort
should be devoted to the rationale section.  If the application is
safety-critical, the bar is higher.  This is currently less structured, but
analogous to, the idea to the Automotive Safety Integrity Levels (ASILs) that
McSCert uses in their automotive hazard analyses.}

\subsection{User Characteristics} \label{SecUserCharacteristics}

\plt{This section summarizes the knowledge/skills expected of the user.
  Measuring usability, which is often a required non-function requirement,
  requires knowledge of a typical user.  As mentioned above, the user is a
  different role from the ``intended reader,'' as given in
  Section~\ref{sec_IntendedReader}.  As in Section~\ref{sec_IntendedReader}, the
  user characteristics should be specific an unambiguous.  For instance, ``The
  end user of \progname{} should have an understanding of undergraduate Level 1
  Calculus and Physics.''}

\subsection{System Constraints}

\begin{itemize}
    \item \progname{} shall operate on consumer smartphones and tablets running
    iOS or Android.
    (C.\refstepcounter{connum}\theconnum\label{C_MobilePlatforms})
    \begin{itemize}
        \item \textbf{Rationale:} Patient families and clinics may use mobile devices running
        either operating system. Supporting both enables preparation
        at home and in clinic without assuming a particular platform.
    \end{itemize}

    \item Third-party assets and libraries included in \progname{} shall have
    licenses allowing their use and redistribution within the application.
    Required copyright notices, licenses, and attributions shall be retained.
    (C.\refstepcounter{connum}\theconnum\label{C_ApplicationLicensing})
    \begin{itemize}
        \item \textbf{Rationale:} The application's source code is licensed under the
        MIT license, while third-party materials have their own licensing.
        These conditions must be satisfied when the application is released.
    \end{itemize}

    \item Third-party materials present in the \progname{} GitHub
    repository shall have licenses allowing redistribution in the form committed to the repository.
    (C.\refstepcounter{connum}\theconnum\label{C_RepositoryLicensing})
    \begin{itemize}
        \item \textbf{Rationale:} Permission to distribute an asset in
        an application does not necessarily mean redistribution of its raw
        files in the project repository is permitted. Repository contents must satisfy
        the applicable licensing conditions (separately from the compiled application).
    \end{itemize}

    \item \progname{} shall be built using a game engine with native support
    for cross-platform deployment to iOS and Android.
    (C.\refstepcounter{connum}\theconnum\label{C_CrossPlatformEngine})
    \begin{itemize}
        \item \textbf{Rationale:} Native deployment support satisfies \conref{C_MobilePlatforms}. The team has
        currently selected the Unity game engine, but the selected engine may
        change provided that it satisfies this constraint.
    \end{itemize}

    \item \progname{} shall provide all included preparation modules and their
    required assets without internet connection after installation.
    (C.\refstepcounter{connum}\theconnum\label{C_OfflineOperation})
    \begin{itemize}
        \item \textbf{Rationale:} Families must be able to prepare at home and
        in clinic, where connectivity may be poor or unavailable.
    \end{itemize}

    \item \progname{} shall operate as a standalone local application and shall
    not initiate network communications during operation, including
    through included third-party libraries.
    (C.\refstepcounter{connum}\theconnum\label{C_NoNetworkCommunications})
    \begin{itemize}
        \item \textbf{Rationale:} Preparation activities do not require external
        services or data transit. Prohibiting network
        communications also prevents included libraries from transmitting usage
        data. Downloading and updating the application through its distribution
        channels are outside normal application operation (i.e., a patient progressing through the application's
        preparation modules).
    \end{itemize}

    \item \progname{} shall not request, collect, or store identifying
    information or clinical records, and shall not integrate with clinic
    information systems.
    (C.\refstepcounter{connum}\theconnum\label{C_PatientPrivacy})
    \begin{itemize}
        \item \textbf{Rationale:} Preparation activities do not require patient
        records. Limiting application data to local gameplay, like character selection and scores,
        supports preparation without requiring identifying information.
    \end{itemize}

    \item The iOS and Android releases of \progname{} shall comply with the Apple App Store and Google Play submission
    requirements, respectively, including applicable policies for applications meant for children, at time of submission.
    (C.\refstepcounter{connum}\theconnum\label{C_StoreCompliance})
    \begin{itemize}
        \item \textbf{Rationale:} These marketplaces are the planned distribution
        channels through which families and clinics will install the application. Compliance
        with their submission requirements is necessary for distribution to the target users.
    \end{itemize}
\end{itemize}

% ============================================================== 
\section{Stakeholder Requirements}

\plt{Technology-independent requirements.  These come from statements made by
the stakeholders.  You will find these through interviews and surveys.  They are
often expressed in natural language.  They are not necessarily testable, but
they should be verifiable.}

\plt{The stakeholder requirements are numbered so that you can have traceability to other parts of your document.}

\noindent \begin{itemize}

\item[SR\refstepcounter{srnum}\thesrnum \label{SR_Label1}] TODO

\item[SR\refstepcounter{srnum}\thesrnum \label{SR_Label2}] TODO

\item[SR\refstepcounter{srnum}\thesrnum \label{SR_Label3}] TODO

\end{itemize}

\plt{You should give your SRs meaningful labels, so that you can refer to them
in other parts of the SRS.  The label names should be meaningful mnemonics, not
numbers.}

% ============================================================== 
\section{User Personas}
\plt{Create 2-5 realistic personas.
  Personas are fictional characters that represent the different user types that
  might use your software.  They are used to help understand the users'
  needs, experiences, behaviors and goals.  They are not real people, but they
  should be based on real data collected from user research.}

\subsection{Persona 1}
\begin{itemize}
\item Name:
\item Description:
\item Goals:
\item Frustrations:
\item Technical Ability:
\end{itemize}

\subsection{Persona 2}
\begin{itemize}
\item Name:
\item Description:
\item Goals:
\item Frustrations:
\item Technical Ability:
\end{itemize}

% ============================================================== 
\section{Use Cases}

The following use cases describe how children and caregivers
interact with \progname{} to prepare for upcoming paediatric
ophthalmology appointments. These use cases apply to both
caregiver-configured and code-based game recommendations.

\refstepcounter{ucnum}
\label{UC_FirstUse}
\subsection{\ucref{UC_FirstUse}: Initial Application Setup}

\textbf{Primary Actor}

Caregiver

\textbf{Preconditions}
\begin{itemize}
    \item The application is installed on a supported device.
    \item The application is being opened for the first time.
\end{itemize}

\textbf{Main Success Scenario}
\begin{enumerate}
    \item The caregiver opens the application.
    \item The system displays an introduction explaining
    the application's purpose.
    \item The system informs the caregiver that the
    application is a preparation aid and does not
    replace professional medical advice.
    \item The caregiver reads and acknowledges the notice.
    \item The system records the acknowledgment locally.
    \item The system allows the caregiver to continue
    to the application's home screen.
\end{enumerate}

\textbf{Alternative Flows}
\begin{itemize}
    \item If the caregiver does not acknowledge the notice,
    the system does not proceed beyond the initial
    acknowledgment step.
    \item If the application is closed during setup,
    the system displays the notice again when the
    application is reopened.
    \item If the caregiver has previously acknowledged
    the notice, the system skips the initial
    acknowledgment step.
\end{itemize}

\textbf{Postconditions}
\begin{itemize}
    \item The caregiver has been informed of the
    application's intended purpose and limitations.
    \item The initial acknowledgment has been
    recorded locally.
    \item The application is ready for further configuration
    or use.
\end{itemize}


\refstepcounter{ucnum}
\label{UC_ProcedureInformation}
\subsection{\ucref{UC_ProcedureInformation}: View Procedure Information}

\textbf{Primary Actor}

Caregiver

\textbf{Preconditions}
\begin{itemize}
    \item The application is open.
    \item Procedure information is available within
    the application.
\end{itemize}

\textbf{Main Success Scenario}
\begin{enumerate}
    \item The caregiver navigates to the procedure
    information or caregiver guide section.
    \item The system displays the available
    ophthalmology procedures.
    \item The caregiver selects a procedure.
    \item The system displays an age-appropriate
    explanation of the selected procedure and
    the preparation game associated with it.
    \item The caregiver reviews the information to
    understand how the game prepares the child
    for the upcoming examination.
    \item The caregiver returns to the previous
    screen or selects another procedure.
\end{enumerate}

\textbf{Alternative Flows}
\begin{itemize}
    \item If the caregiver selects a different procedure,
    the system displays the information associated
    with the newly selected procedure.
    \item If the caregiver exits the guide,
    the system returns to the previous screen
    without changing game progress.
\end{itemize}

\textbf{Postconditions}
\begin{itemize}
    \item The caregiver has access to information
    explaining the selected procedure and game.
    \item No changes are made to the child's
    game progress.
\end{itemize}


\refstepcounter{ucnum}
\label{UC_SelectGame}
\subsection{\ucref{UC_SelectGame}: Select a Preparation Game}

\textbf{Primary Actor}

Child

\textbf{Preconditions}
\begin{itemize}
    \item The application is open.
    \item At least one preparation game is available
    through the current recommendation selection.
\end{itemize}

\textbf{Main Success Scenario}
\begin{enumerate}
    \item The child navigates to the game selection screen.
    \item The system displays the preparation games
    available under the current recommendation.
    \item The child selects a game using the
    displayed visual controls.
    \item The system presents the selected game's
    introduction and instructions.
    \item The child begins the game.
\end{enumerate}

\textbf{Alternative Flows}
\begin{itemize}
    \item If no games have been selected, the system
    informs the caregiver that game configuration
    is required.
    \item If the child returns to the home screen
    before starting the game, the system exits
    the selection process without updating progress.
    \item If the selected game has existing progress,
    the system uses the game's supported
    continuation or restart behaviour.
\end{itemize}

\textbf{Postconditions}
\begin{itemize}
    \item The selected game has been opened.
    \item The child can begin an age-appropriate
    preparation experience.
    \item Progress for unrelated games remains unchanged.
\end{itemize}


\refstepcounter{ucnum}
\label{UC_VisionGame}
\subsection{\ucref{UC_VisionGame}: Complete Vision Testing Preparation}

\textbf{Primary Actor}

Child

\textbf{Preconditions}
\begin{itemize}
    \item The vision testing game is available to the child.
    \item The child has selected the vision testing game.
\end{itemize}

\textbf{Main Success Scenario}
\begin{enumerate}
    \item The system introduces the vision testing procedure
    through a child-friendly animation demonstrating how
    an eye examination may be performed.

    \item The system explains that the child will practice
    identifying and matching different symbols displayed
    on the screen.

    \item The system displays a reference symbol and a
    set of possible matching symbols.

    \item The child selects the symbol that matches
    the reference symbol or chooses to skip the round.

    \item The system provides encouraging feedback
    based on the child's selection or acknowledges
    the decision to skip.

    \item Steps 3--5 repeat with different symbols
    until all symbol-matching rounds have been
    attempted or skipped.

    \item The system displays a completion message
    and records the game as completed.
\end{enumerate}

\textbf{Alternative Flows}
\begin{itemize}
    \item If the child selects an incorrect symbol,
    the system provides non-punitive feedback
    and allows another attempt.

    \item If the child has difficulty identifying
    a symbol, the child may skip the current round.
    The system proceeds to the next round without
    requiring a correct answer.

    \item If the child requests assistance, a caregiver
    may help explain the instructions.

    \item If the child exits the game before reaching
    the end of all rounds, the system returns to
    the home screen without marking the game
    as completed.
\end{itemize}

\textbf{Postconditions}
\begin{itemize}
    \item The child has been introduced to symbol
    recognition activities associated with vision testing.

    \item If all rounds have been attempted or skipped,
    the game's completion status is updated locally.

    \item The application does not assess visual acuity
    or generate diagnostic results.
\end{itemize}


\refstepcounter{ucnum}
\label{UC_EyeDropGame}
\subsection{\ucref{UC_EyeDropGame}: Complete Eye-Drop Preparation}

\textbf{Primary Actor}

Child

\textbf{Preconditions}
\begin{itemize}
    \item The eye-drop preparation game is available
    to the child.
    \item The child has selected the eye-drop game.
\end{itemize}

\textbf{Main Success Scenario}
\begin{enumerate}
    \item The system introduces a virtual character
    preparing to receive eye drops during an
    ophthalmology appointment.

    \item The system explains through child-friendly
    narration that a clinical professional may
    administer eye drops during an examination.

    \item The system displays a virtual character
    and demonstrates a preparation action,
    such as tilting the character's head back.

    \item The child interacts with the virtual
    character to perform the demonstrated
    preparation action or chooses to skip it.

    \item The system provides encouraging feedback
    or acknowledges the decision to skip.

    \item Steps 3--5 repeat for the remaining
    preparation actions until all have been
    attempted or skipped.

    \item The system displays an animation of
    a clinical professional administering
    eye drops to the virtual character.

    \item The system explains that receiving eye drops
    may feel unfamiliar and that the child can
    communicate discomfort to the clinician.

    \item The child watches the simulation or chooses
    to skip to the end of the demonstration.

    \item The system displays a completion message
    and records the game as completed.
\end{enumerate}

\textbf{Alternative Flows}
\begin{itemize}
    \item If the child performs a preparation action
    incorrectly, the system provides a visual hint
    and allows another attempt.

    \item If the child cannot complete a preparation
    action, the child may skip it and continue
    to the next stage of the simulation.

    \item If the child is uncomfortable watching
    the eye-drop administration animation,
    the child may skip the demonstration.

    \item If the child exits the game before reaching
    the end of the simulation, the system returns
    to the home screen without marking the
    game as completed.
\end{itemize}

\textbf{Postconditions}
\begin{itemize}
    \item The child has been introduced to the
    preparation steps and experience associated
    with receiving eye drops.

    \item If all stages have been attempted or skipped,
    the game's completion status is updated locally.

    \item The application does not instruct the child
    to administer actual medication or independently
    perform medical procedures.
\end{itemize}


\refstepcounter{ucnum}
\label{UC_SlitLampGame}
\subsection{\ucref{UC_SlitLampGame}: Complete Slit-Lamp Preparation}

\textbf{Primary Actor}

Child

\textbf{Preconditions}
\begin{itemize}
    \item The slit-lamp preparation game is available
    to the child.
    \item The child has selected the slit-lamp game.
\end{itemize}

\textbf{Main Success Scenario}
\begin{enumerate}
    \item The system introduces the slit-lamp examination
    through a child-friendly animation of a virtual
    examination room.

    \item The system displays a simulated slit-lamp
    device and explains how a clinician uses it
    to examine the child's eyes.

    \item The system demonstrates a positioning action,
    such as placing the virtual character's chin
    on the examination device's chin rest.

    \item The child uses the displayed controls
    to perform the demonstrated positioning
    action or chooses to skip the exercise.

    \item The system provides visual feedback
    indicating whether the character is positioned
    correctly or acknowledges the decision to skip.

    \item Steps 3--5 repeat for the remaining
    positioning exercises, including forehead
    placement, until all exercises have been
    attempted or skipped.

    \item The system introduces a short simulation
    demonstrating the importance of remaining
    relatively still during the examination.

    \item The child participates in the simulated
    positioning exercise or chooses to skip it.

    \item The system provides encouraging feedback
    and displays an animation representing
    the completion of the examination.

    \item The system displays a completion message
    and records the game as completed.
\end{enumerate}

\textbf{Alternative Flows}
\begin{itemize}
    \item If the virtual character is positioned
    incorrectly, the system provides visual
    guidance and allows the child to adjust
    its position.

    \item If the child has difficulty with a
    positioning exercise, the child may skip
    it without being required to achieve
    the correct position.

    \item If the child cannot complete the simulated
    stillness exercise, the system allows
    the child to skip it and continue
    to the final demonstration.

    \item If the child exits the game before
    reaching its final stage, the system returns
    to the home screen without marking
    the game as completed.
\end{itemize}

\textbf{Postconditions}
\begin{itemize}
    \item The child has been introduced to the
    slit-lamp examination device and the
    positioning involved in the procedure.

    \item If all exercises have been attempted
    or skipped, the game's completion status
    is updated locally.

    \item The application does not produce
    clinical measurements or diagnostic results.
\end{itemize}


\refstepcounter{ucnum}
\label{UC_ExitGame}
\subsection{\ucref{UC_ExitGame}: Exit a Preparation Game Before Completion}

\textbf{Primary Actor}

Child

\textbf{Preconditions}
\begin{itemize}
    \item The child is currently playing a preparation game,
    such as the vision testing, eye-drop, or slit-lamp game.
    \item The child has not yet reached the end of the game.
\end{itemize}

\textbf{Main Success Scenario}
\begin{enumerate}
    \item While participating in a preparation game, the
    child decides to stop playing before completing
    all the game's rounds or exercises.

    \item The child selects the home button displayed
    on the game screen.

    \item The system stops the current game interaction,
    including any ongoing animations, narration,
    and game specific input processing.

    \item The system closes the game screen and returns
    the child to the game selection screen.

    \item The system displays the available preparation
    games, including the game the child just exited.

    \item The system retains the completion records
    of previously completed games and does not mark
    the interrupted game as completed.

    \item The child may select another preparation
    game or choose to replay the game they exited.
\end{enumerate}

\textbf{Alternative Flows}
\begin{itemize}
    \item \textbf{Exit During a Game Animation:}
    If the child selects the home button while an
    animation or narration is playing, the system
    stops the animation and audio before returning
    to the game selection screen.

    \item \textbf{Exit After Skipping Multiple Rounds:}
    If the child has skipped one or more rounds
    but has not reached the end of the game,
    the system allows the child to exit without
    requiring any skipped rounds to be completed.

    \item \textbf{Application Interruption:}
    If the application is closed or interrupted
    during gameplay, the system retains previously
    completed game records. When the application
    is reopened, the child can access the game
    selection screen and restart the interrupted game.

    \item \textbf{Replay an Interrupted Game:}
    If the child selects the previously interrupted
    game again, the system starts the game from
    its introduction rather than automatically
    recording the previous attempt as completed.
\end{itemize}

\textbf{Postconditions}
\begin{itemize}
    \item The child is returned to the game selection
    screen, where the available preparation games
    are displayed.

    \item The interrupted game's animations, audio,
    and interactive controls are no longer running.

    \item The interrupted game is not recorded as
    completed unless its completion conditions
    were already satisfied.

    \item Completion records for previously completed
    games remain unchanged.

    \item The child can start another game or
    replay the interrupted game without
    completing the previous attempt.
\end{itemize}


\refstepcounter{ucnum}
\label{UC_CompleteGame}
\subsection{\ucref{UC_CompleteGame}: Receive Game Completion Feedback}

\textbf{Primary Actor}

Child

\textbf{Preconditions}
\begin{itemize}
    \item A preparation game is active.
    \item The child has reached the game's
    defined completion conditions.
\end{itemize}

\textbf{Main Success Scenario}
\begin{enumerate}
    \item The child completes the required
    interactions within the game.
    \item The system recognizes that
    the game's completion conditions
    have been satisfied.
    \item The system records the game
    as completed.
    \item The system presents positive
    completion feedback or a reward.
    \item The child returns to the
    game selection screen.
    \item The system displays the game
    as completed.
\end{enumerate}

\textbf{Alternative Flows}
\begin{itemize}
    \item If the game is exited before
    completion, the system does not
    issue a completion reward.
    \item If the child replays a previously
    completed game, the system allows
    the replay without removing its
    completed status.
\end{itemize}

\textbf{Postconditions}
\begin{itemize}
    \item The game's completion status
    is stored locally.
    \item The child receives positive
    feedback for completing the game.
    \item Completion does not indicate
    clinical readiness or guarantee
    cooperation during an appointment.
\end{itemize}


\refstepcounter{ucnum}
\label{UC_ViewProgress}
\subsection{\ucref{UC_ViewProgress}: View Preparation Progress}

\textbf{Primary Actor}

Caregiver

\textbf{Preconditions}
\begin{itemize}
    \item The application is open.
    \item The caregiver has configured or entered a
    recommendation for the child's preparation games.
    \item The system has access to the locally stored
    completion status of each game.
\end{itemize}

\textbf{Main Success Scenario}
\begin{enumerate}
    \item The caregiver selects the progress button
    from the application's home screen.

    \item The system displays the preparation progress
    screen, showing the games currently recommended
    for the child.

    \item The system retrieves the locally stored
    completion status of each recommended game.

    \item For each game, the system displays its name
    and one of the following completion statuses:
    \begin{itemize}
        \item \textbf{Not Started:} The child has not
        started the game.
        \item \textbf{In Progress:} The child has
        started the game but exited before reaching
        its final stage.
        \item \textbf{Completed:} The child has
        reached the end of the game after attempting
        or skipping its activities.
    \end{itemize}

    \item The system visually distinguishes completed
    games from incomplete games using status
    indicators that do not rely solely on colour.

    \item Steps 3--5 are repeated for each game in
    the current recommendation until the progress
    information for all recommended games is displayed.

    \item The caregiver reviews the status of each
    game to identify which preparation games the
    child has completed and which remain unfinished.

    \item The caregiver selects the home button
    to return to the application's home screen.
\end{enumerate}

\textbf{Alternative Flows}
\begin{itemize}
    \item \textbf{No Games Started:}
    If the child has not started any recommended
    games, the system displays all recommended
    games with the status \textit{Not Started}.

    \item \textbf{Game Exited Before Completion:}
    If the child previously exited a game before
    reaching its final stage, the system displays
    that game with the status \textit{In Progress}.
    The game is not marked as completed.

    \item \textbf{Game Completed with Skipped Rounds:}
    If the child reached the end of a game after
    attempting or skipping its rounds, the system
    displays the game with the status
    \textit{Completed}.

    \item \textbf{No Games Configured:} CHANGE-REQ
    If no games have been selected through caregiver
    configuration or a recommendation code, the
    system informs the caregiver that there are
    no recommended games to display.

    \item \textbf{Progress Previously Cleared:}
    If the caregiver has deleted the locally stored
    game progress, the system displays each
    recommended game with the status
    \textit{Not Started}.
\end{itemize}

\textbf{Postconditions}
\begin{itemize}
    \item The caregiver has viewed the completion
    status of each currently recommended game.

    \item The progress screen has not modified
    any stored game completion records.

    \item All game progress remains stored locally
    on the device.

    \item The displayed completion statuses represent
    participation in the preparation games and
    do not indicate clinical performance, visual
    ability, or readiness for an examination.
\end{itemize}


\refstepcounter{ucnum}
\label{UC_AdjustPreferences}
\subsection{\ucref{UC_AdjustPreferences}: Change Application Preferences}

\textbf{Primary Actor}

Caregiver

\textbf{Preconditions}
\begin{itemize}
    \item The application is open.
    \item The caregiver can access the application settings.
    \item A nickname has previously been saved.
\end{itemize}

\textbf{Main Success Scenario}
\begin{enumerate}
    \item The caregiver selects the settings button
    from the application's home screen.

    \item The system displays the available settings,
    including the current nickname and display
    and audio preferences.

    \item The caregiver selects the nickname field,
    which displays the currently saved nickname.

    \item The system allows the caregiver to edit
    the existing nickname.

    \item The caregiver replaces the existing nickname
    with a new nickname of their choice and
    selects the save button.

    \item The system saves the updated nickname
    locally on the device.

    \item The system updates the nickname field
    to display the newly saved nickname.

    \item The caregiver selects the home button
    to return to the application's home screen.

    \item The system displays the updated nickname
    wherever it is used within the application.
\end{enumerate}

\textbf{Alternative Flows}
\begin{itemize}
    \item \textbf{Empty Nickname:}
    If the caregiver clears the nickname field,
    the system allows it to remain empty and
    uses default text where personalization
    would normally appear.

    \item \textbf{Cancel Changes:}
    If the caregiver cancels the changes before
    saving, the system retains the previously
    saved nickname.

    \item \textbf{Modify Other Preferences:}
    If the caregiver selects an audio or display
    setting instead, the system allows the
    caregiver to modify and save that preference.

    \item \textbf{Application Restart:}
    If the application is closed and reopened
    after saving, the system retrieves and
    displays the updated nickname.
\end{itemize}

\textbf{Postconditions}
\begin{itemize}
    \item The previously saved nickname has
    been replaced with the caregiver's
    newly chosen nickname.

    \item The updated nickname is displayed
    wherever personalization is supported
    within the application.

    \item The updated nickname remains available
    after the application is closed and reopened.

    \item Existing game recommendations and
    completion records remain unchanged.

    \item No nickname information is transmitted
    to an external server.
\end{itemize}


\refstepcounter{ucnum}
\label{UC_DeleteData}
\subsection{\ucref{UC_DeleteData}: Clear Locally Stored User Data}

\textbf{Primary Actor}

Caregiver

\textbf{Preconditions}
\begin{itemize}
    \item The application is installed.
    \item The caregiver can access
    the data management settings.
\end{itemize}

\textbf{Main Success Scenario}
\begin{enumerate}
    \item The caregiver opens
    the data management section.
    \item The system displays
    an option to clear
    locally stored user data.
    \item The caregiver selects
    the option to clear data.
    \item The system requests
    confirmation of the action.
    \item The caregiver confirms
    the deletion.
    \item The system clears locally
    stored preferences, game
    selections or recommendations,
    and game progress.
    \item The system returns
    to the default setup state.
\end{enumerate}

\textbf{Alternative Flows}
\begin{itemize}
    \item If the caregiver cancels
    the confirmation, no data
    is deleted.
    \item If there is no user
    data to clear, the system
    indicates that no stored
    progress or preferences
    are available for deletion.
\end{itemize}

\textbf{Postconditions}
\begin{itemize}
    \item The selected locally
    stored user data has
    been cleared.
    \item The application remains
    installed and usable.
    \item The application's
    built-in games and educational
    content remain available.
\end{itemize}


\refstepcounter{ucnum}
\label{UC_OfflineUse}
\subsection{\ucref{UC_OfflineUse}: Access Preparation Games Offline}

\textbf{Primary Actor}

Child

\textbf{Preconditions}
\begin{itemize}
    \item The application has been installed successfully
    on a supported device.

    \item CHANGE-REQ The caregiver has previously configured the
    recommended games or entered a valid recommendation
    code, and the selection has been saved locally.

    \item The application is closed.

    \item The device has no active Internet connection.
\end{itemize}

\textbf{Main Success Scenario}
\begin{enumerate}
    \item The child opens the application while the
    device has no Internet connection.

    \item The system loads the previously saved game
    recommendations, preferences, and progress
    from the device's local storage.

    \item The system displays the application's
    home screen without requiring an Internet
    connection.

    \item The child selects the button to view
    the available preparation games.

    \item The system displays the previously
    recommended games, such as the Vision Testing,
    Eye-Drop, and Slit-Lamp Preparation Games.

    \item The child selects the Vision Testing Game.

    \item The system opens the game and presents
    its introduction and instructions using
    locally available content.

    \item The system displays a reference symbol
    and a set of possible matching symbols.

    \item The child selects the matching symbol
    or chooses to skip the current round.

    \item The system provides encouraging feedback
    or acknowledges the decision to skip.

    \item Steps 8--10 repeat with different symbols
    until all required rounds have been attempted
    or skipped.

    \item The system displays a completion message
    and records the Vision Testing Game as completed
    in the device's local storage.

    \item The child selects the home button to
    return to the application's home screen.

    \item The system displays the updated game
    completion status without requiring an
    Internet connection.
\end{enumerate}

\textbf{Alternative Flows}
\begin{itemize}

    \item \textbf{Select a Different Game:}
    If the child selects the Eye-Drop or Slit-Lamp
    Preparation Game instead of the Vision Testing
    Game in Step 6, the system launches the selected
    game using locally available content. The child
    completes or skips its individual exercises,
    and the system saves its completion status locally.

    \item \textbf{Exit Before Completion:}
    If the child selects the home button before
    reaching the end of the game, the system stops
    the game's animations, audio, and interactions
    and returns to the home screen. The interrupted
    game is not marked as completed.

    \item \textbf{Internet Connection Lost During Gameplay:}
    If the device initially has an Internet connection
    but loses it during gameplay, the current game
    continues without interruption. The child can
    finish the game, and its completion status is
    saved locally.

    \item \textbf{Previously Completed Game:}
    If the child selects a game that was previously
    completed, the system allows the game to be
    replayed without resetting its completed status.

\end{itemize}

\textbf{Postconditions}
\begin{itemize}
    \item The child has accessed and played a
    recommended preparation game without requiring
    an Internet connection.

    \item If the child reached the end of the game,
    its completion status is updated and stored
    locally on the device.

    \item Previously saved recommendations,
    preferences, and game progress remain available.

    \item No external server or network connection
    was required to load the game, process interactions,
    or save completion progress.
\end{itemize}




% Option 2 Use case Scenarios
CHANGE-REQ
\refstepcounter{ucnum}
\label{UC_ConfigureGames}
\subsection{\ucref{UC_ConfigureGames}: Configure Recommended Games}

\textbf{Primary Actor}

Caregiver

\textbf{Preconditions}
\begin{itemize}
    \item The application is installed and open.
    \item The caregiver has received recommendations
    from the child's clinician regarding the
    procedures expected during the appointment.
    \item The caregiver can access the
    game configuration settings.
\end{itemize}

\textbf{Main Success Scenario}
\begin{enumerate}
    \item The caregiver navigates to
    the game configuration section.
    \item The system displays the available
    procedure games and their descriptions.
    \item The caregiver selects the age group
    appropriate for the child.
    \item The caregiver selects the games
    recommended by the clinician.
    \item The system displays the selected
    games for review.
    \item The caregiver confirms the selection.
    \item The system saves the selected age
    group and enabled games locally.
    \item The system displays the selected
    games in the child's game selection screen.
\end{enumerate}

\textbf{Alternative Flows}
\begin{itemize}
    \item If the caregiver makes no game selections,
    the system informs the caregiver that
    at least one game must be selected
    before completing the configuration.
    \item If the caregiver cancels the configuration,
    the system retains the previously saved settings.
    \item If the caregiver selects an unsupported
    age group for a game, the system indicates
    that the game is not available for that
    age group.
\end{itemize}

\textbf{Postconditions}
\begin{itemize}
    \item The caregiver's game selections
    and age-group preference are stored locally.
    \item The child can access the games
    selected by the caregiver.
    \item No patient account or clinical
    database connection is required.
\end{itemize}


\refstepcounter{ucnum}
\label{UC_ModifyGames}
\subsection{\ucref{UC_ModifyGames}: Modify Recommended Games}

\textbf{Primary Actor}

Caregiver

\textbf{Preconditions}
\begin{itemize}
    \item The application has an existing
    caregiver configuration.
    \item The caregiver can access the
    game configuration settings.
\end{itemize}

\textbf{Main Success Scenario}
\begin{enumerate}
    \item The caregiver opens the
    game configuration section.
    \item The system displays the currently
    enabled games.
    \item The caregiver reviews the
    existing game selections.
    \item The caregiver enables or disables
    games based on updated clinical
    recommendations or preparation needs.
    \item The system displays the revised
    selection for confirmation.
    \item The caregiver confirms the changes.
    \item The system saves the updated
    configuration locally.
    \item The system updates the games
    available on the child's game
    selection screen.
\end{enumerate}

\textbf{Alternative Flows}
\begin{itemize}
    \item If the caregiver cancels the changes,
    the previous configuration remains active.
    \item If all games are disabled, the
    system informs the caregiver that
    no preparation games are currently selected.
    \item If a previously completed game
    is disabled, its completion record
    is retained unless the caregiver
    separately clears the stored data.
\end{itemize}

\textbf{Postconditions}
\begin{itemize}
    \item The updated game configuration
    is stored locally.
    \item The child's game selection
    screen reflects the updated selections.
    \item Existing game completion records
    remain unchanged.
\end{itemize}


% Option 3 Use case scenarios
% \refstepcounter{ucnum}
% \label{UC_EnterCode}
% \subsection{\ucref{UC_EnterCode}: Enter a Procedure Recommendation Code}

% \textbf{Primary Actor}

% Caregiver

% \textbf{Preconditions}
% \begin{itemize}
%     \item The application is installed and open.
%     \item The caregiver has received a recommendation
%     code from the child's clinician.
%     \item The application contains a locally
%     available mapping of supported
%     recommendation codes to games.
% \end{itemize}

% \textbf{Main Success Scenario}
% \begin{enumerate}
%     \item The caregiver navigates to
%     the recommendation code entry screen.
%     \item The system prompts the caregiver
%     to enter a recommendation code.
%     \item The caregiver enters the code
%     provided by the clinician.
%     \item The system validates the code
%     against its local recommendation catalogue.
%     \item The system identifies the age
%     group associated with the code.
%     \item The system identifies the
%     procedure games associated with the code.
%     \item The system displays the age group
%     and recommended games for review.
%     \item The caregiver confirms the recommendation.
%     \item The system stores the active
%     recommendation locally.
%     \item The system displays the corresponding
%     games on the child's game selection screen.
% \end{enumerate}

% \textbf{Alternative Flows}
% \begin{itemize}
%     \item If the caregiver enters an invalid
%     or unrecognized code, the system displays
%     an error message and allows another attempt.
%     \item If the caregiver cancels the operation,
%     the system retains the previously
%     active recommendation, if any.
%     \item If the caregiver does not confirm
%     the displayed recommendation, the system
%     does not apply the new code.
%     \item If the device has no Internet
%     connection, the system validates
%     the code using its locally
%     available catalogue.
% \end{itemize}

% \textbf{Postconditions}
% \begin{itemize}
%     \item A valid recommendation code
%     has been applied to the application.
%     \item The appropriate age group
%     and games have been selected.
%     \item The active recommendation
%     is stored locally.
%     \item No patient account or external
%     clinical system is required.
% \end{itemize}


% \refstepcounter{ucnum}
% \label{UC_ChangeCode}
% \subsection{\ucref{UC_ChangeCode}: Change the Active Recommendation Code}

% \textbf{Primary Actor}

% Caregiver

% \textbf{Preconditions}
% \begin{itemize}
%     \item The application has an active
%     recommendation code.
%     \item The caregiver has received
%     a different recommendation code
%     or wishes to update the current selection.
% \end{itemize}

% \textbf{Main Success Scenario}
% \begin{enumerate}
%     \item The caregiver opens the
%     recommendation settings.
%     \item The system displays the currently
%     active recommendation and its
%     associated games.
%     \item The caregiver selects the option
%     to change the recommendation code.
%     \item The system prompts the caregiver
%     to enter a new code.
%     \item The caregiver enters the new code.
%     \item The system validates the code
%     using its local recommendation catalogue.
%     \item The system displays the new age
%     group and associated procedure games.
%     \item The caregiver confirms the changes.
%     \item The system replaces the active
%     recommendation with the new selection.
%     \item The system updates the child's
%     game selection screen.
% \end{enumerate}

% \textbf{Alternative Flows}
% \begin{itemize}
%     \item If the new code is invalid,
%     the system displays an error message
%     and retains the existing recommendation.
%     \item If the caregiver cancels the operation,
%     the previously active recommendation
%     remains unchanged.
%     \item If the new recommendation contains
%     games that were previously completed,
%     their stored completion records are retained.
%     \item If the caregiver clears the active
%     recommendation without entering a new code,
%     the system returns to an unconfigured state.
% \end{itemize}

% \textbf{Postconditions}
% \begin{itemize}
%     \item The new recommendation code
%     replaces the previous active code
%     if the caregiver confirmed the change.
%     \item The child's available games and
%     age-appropriate content reflect
%     the new recommendation.
%     \item Previously stored game progress
%     remains unchanged.
%     \item The updated recommendation
%     is stored locally.
% \end{itemize}

% ============================================================== 
\section{Functional Requirements}
\plt{Functional requirements describe what the system should do.  They are typically
expressed as a set of conditions or capabilities that the system must have.  These
requirements are usually derived from the stakeholder requirements and use cases.}

\plt{Use SHALL statements. Requirements must be atomic, testable, and traceable.}

\noindent \begin{itemize}

\item[FR\refstepcounter{reqnum}\thereqnum \label{FR_Label1}:] The system shall ...

\item[FR\refstepcounter{reqnum}\thereqnum \label{FR_Label2}:] The system shall ...

\item[FR\refstepcounter{reqnum}\thereqnum \label{FR_Label3}:] The system shall ...

\end{itemize}

\plt{The requirements are given labels so that they can be referenced in other
parts of the SRS.  The label names should be meaningful mnemonics, not numbers.}

% ============================================================== 
\section{Data Requirements}

This section describes the information managed by \progname{}
to support age appropriate preparation activities, caregiver
configuration, and local activity progress. The application
does not require patient accounts or integration with clinical
information systems. All application data is maintained locally
on the user's device.

\subsection{Major Data Entities}

The system manages the following major data entities:

\begin{itemize}

    \item \textbf{Procedure:}
    Represents a common paediatric ophthalmology procedure
    introduced through the application, such as vision testing,
    eye drop administration, or slit lamp examinations.
    Each procedure is associated with one or more preparation
    activities.

    \item \textbf{Preparation Activity:}
    Represents an interactive game or educational activity
    designed to familiarize children with a particular
    ophthalmology procedure. Each activity is associated
    with a procedure and an intended age group.

    \item \textbf{Caregiver Configuration:}
    Represents the selection of preparation activities
    enabled by a parent or caregiver based on the
    recommendations provided by the child's clinician.
    The configuration determines which activities appear
    in the child's preparation experience.

    \item \textbf{Local Preferences:}
    Represents optional settings associated with the
    application, including the child's nickname, age
    group, and display or accessibility preferences.
    These settings are stored locally and do not
    require a patient account.

    \item \textbf{Activity Progress:}
    Represents the completion status of preparation
    activities, allowing children and caregivers to
    view which activities have been completed.
    This information reflects application progress
    rather than clinical performance or assessment results.

\end{itemize}

The relationship between these entities is as follows:
each Procedure may have multiple Preparation Activities.
A Caregiver Configuration identifies the activities
selected for the child. Activity Progress records the
completion status of individual activities, while
Local Preferences determine how the application is
presented to the user.

\subsection{Data Dictionary}

The following data dictionary defines the major
attributes associated with each entity.

\begin{longtable}{
    p{0.16\textwidth}
    p{0.24\textwidth}
    p{0.16\textwidth}
    p{0.33\textwidth}
}

\caption{Data Dictionary for Caregiver Configured Preparation}
\label{tab:DataDictionary}\\

\toprule
\textbf{Entity} &
\textbf{Attribute} &
\textbf{Type} &
\textbf{Description} \\
\midrule
\endfirsthead

\toprule
\textbf{Entity} &
\textbf{Attribute} &
\textbf{Type} &
\textbf{Description} \\
\midrule
\endhead

Procedure &
procedure\_id &
String &
Unique identifier for an ophthalmology procedure. \\
\midrule

Procedure &
procedure\_name &
String &
Human readable name of the procedure. \\
\midrule

Procedure &
procedure\_description &
String &
Clinically reviewed educational description of the
procedure. \\
\midrule

Preparation Activity &
activity\_id &
String &
Unique identifier for a preparation activity. \\
\midrule

Preparation Activity &
procedure\_id &
String &
Identifier of the procedure associated with the
activity. \\
\midrule

Preparation Activity &
activity\_name &
String &
Name displayed for the preparation activity. \\
\midrule

Preparation Activity &
age\_groups &
Set of Strings &
Age groups for which the activity is intended. \\
\midrule

Caregiver Configuration &
enabled\_activities &
Set of Strings &
Identifiers of the preparation activities selected
by the caregiver. \\
\midrule

Caregiver Configuration &
configuration\_status &
Enumeration &
Indicates whether the caregiver has configured
the activity selection. Values: Unconfigured,
Configured. \\
\midrule

Local Preferences &
nickname &
Optional String &
Optional nickname used to personalize the
application without requiring a legal name. \\
\midrule

Local Preferences &
age\_group &
Optional String &
Age group selected by the caregiver to determine
which activity versions are appropriate. \\
\midrule

Local Preferences &
display\_settings &
Settings Record &
Locally stored display, audio, and accessibility
preferences. \\
\midrule

Activity Progress &
activity\_id &
String &
Identifier of the activity associated with the
progress record. \\
\midrule

Activity Progress &
completion\_status &
Enumeration &
Current activity status. Values: Not Started,
In Progress, Completed. \\

\bottomrule

\end{longtable}

\subsection{Privacy and Data Retention}

\progname{} follows a data minimization approach, collecting
and retaining only the information necessary to support
caregiver configured preparation activities and local progress
tracking. The application's privacy and data management
practices are outlined below:

\begin{itemize}

    \item \textbf{Data Collection:}
    The application does not require patient accounts,
    medical records, appointment details, or personally
    identifying information. Optional nicknames are used
    solely for personalization and do not need to
    correspond to the child's legal name.

    \item \textbf{Caregiver Configuration:}
    Caregivers select preparation activities based on
    recommendations provided by the child's clinician.
    The application does not access clinical information
    systems or independently retrieve appointment specific
    medical information.

    \item \textbf{Data Storage and Protection:}
    User preferences, caregiver configurations, and activity
    progress shall be stored locally on the user's device
    and protected from unauthorized access by other applications.
    The application shall not transmit this information
    to external servers or require user accounts.

    \item \textbf{Data Retention and Deletion:}
    Locally stored information is retained until
    deleted by the caregiver or through application
    removal, subject to the device's backup and
    restoration behaviour. Caregivers can clear
    stored preferences, activity selections, and
    progress through the application's data
    management functionality.

    \item \textbf{Consent and Access:}
    The application does not require registration
    or the collection of personal information to
    access its preparation activities. Caregivers
    control the configuration of activities
    presented to the child.

    \item \textbf{Research and Regulatory Compliance:}
    Any personal information collected separately
    for research or usability evaluation is outside
    the application's operational data storage
    and is subject to applicable institutional
    ethics and privacy requirements.

\end{itemize}


% \section{Data Requirements}

% This section describes the information managed by \progname{}
% to support procedure recommendations through predefined
% codes. The application does not require patient accounts
% or access to clinical information systems.

% Clinicians may provide caregivers with a generic
% recommendation code corresponding to an age group and
% a set of preparation activities. The caregiver enters
% the code into the application, which identifies the
% associated activities using a locally available mapping.

% The recommendation codes described in this section
% are an initial draft and will be reviewed with the project
% supervisors before finalization.

% \subsection{Major Data Entities}

% The system manages the following major data entities:

% \begin{itemize}

%     \item \textbf{Procedure:}
%     Represents a common paediatric ophthalmology
%     procedure introduced through the application,
%     such as vision testing, eye drop administration,
%     or slit lamp examinations.

%     \item \textbf{Preparation Activity:}
%     Represents an interactive game or educational
%     activity associated with an ophthalmology
%     procedure. Activities may have different
%     versions designed for particular age groups.

%     \item \textbf{Recommendation Code:}
%     Represents a predefined mapping between a
%     generic recommendation code, an intended
%     age group, and a set of preparation activities.
%     These codes are not patient-specific and do
%     not contain personal or appointment information.

%     \item \textbf{Local Preferences:}
%     Represents the locally maintained application
%     settings, including an optional nickname,
%     display preferences, and the active
%     recommendation selection.

%     \item \textbf{Activity Progress:}
%     Represents the completion status of individual
%     preparation activities. Progress is maintained
%     locally and is not treated as a clinical
%     assessment or measure of visual ability.

% \end{itemize}

% Each Procedure may be associated with multiple
% Preparation Activities. Each Recommendation Code
% identifies an age group and a predefined set of
% activities associated with one or more procedures.

% When a valid code is entered, the application
% retrieves the corresponding activity selection
% from its local recommendation catalogue.
% The active recommendation determines which
% activities are presented to the child.

% Activity Progress is maintained independently
% for each preparation activity.

% \subsection{Data Dictionary}

% The following data dictionary defines the major
% attributes associated with each entity.

% \begin{longtable}{
%     p{0.16\textwidth}
%     p{0.22\textwidth}
%     p{0.16\textwidth}
%     p{0.31\textwidth}
% }

% \caption{Data Dictionary for Code Based Preparation}
% \label{tab:DataDictionary}\\

% \toprule
% \textbf{Entity} &
% \textbf{Attribute} &
% \textbf{Type} &
% \textbf{Description} \\
% \midrule
% \endfirsthead

% \toprule
% \textbf{Entity} &
% \textbf{Attribute} &
% \textbf{Type} &
% \textbf{Description} \\
% \midrule
% \endhead

% Procedure &
% procedure\_id &
% String &
% Unique identifier for an ophthalmology procedure. \\
% \midrule

% Procedure &
% procedure\_name &
% String &
% Human-readable name of the procedure. \\
% \midrule

% Procedure &
% procedure\_description &
% String &
% Clinically reviewed educational description
% of the procedure. \\
% \midrule

% Preparation Activity &
% activity\_id &
% String &
% Unique identifier for a preparation activity. \\
% \midrule

% Preparation Activity &
% procedure\_id &
% String &
% Identifier of the associated ophthalmology procedure. \\
% \midrule

% Preparation Activity &
% activity\_name &
% String &
% Name displayed for the preparation activity. \\
% \midrule

% Preparation Activity &
% age\_groups &
% Set of Strings &
% Age groups for which the activity is intended. \\
% \midrule

% Recommendation Code &
% recommendation\_code &
% String &
% Unique generic code identifying an age group
% and preparation activity selection. \\
% \midrule

% Recommendation Code &
% age\_group &
% Enumeration &
% Age band associated with the code.
% Provisional values: Ages 3--5 and Ages 5--7. \\
% \midrule

% Recommendation Code &
% procedure\_ids &
% Set of Strings &
% Identifiers of the procedures included
% in the recommendation. \\
% \midrule

% Recommendation Code &
% activity\_ids &
% Set of Strings &
% Identifiers of the preparation activities
% associated with the recommendation. \\
% \midrule

% Recommendation Code &
% recommendation\_description &
% String &
% Human-readable description of the
% recommended preparation activities. \\
% \midrule

% Local Preferences &
% nickname &
% Optional String &
% Optional nickname used for local personalization. \\
% \midrule

% Local Preferences &
% active\_recommendation &
% Optional String &
% Identifier of the currently selected
% recommendation code. \\
% \midrule

% Local Preferences &
% display\_settings &
% Settings Record &
% Locally stored display, audio, and
% accessibility preferences. \\
% \midrule

% Activity Progress &
% activity\_id &
% String &
% Identifier of the activity associated
% with the progress record. \\
% \midrule

% Activity Progress &
% completion\_status &
% Enumeration &
% Current activity status. Values:
% Not Started, In Progress, Completed. \\

% \bottomrule

% \end{longtable}

% \subsubsection{Initial Draft of Recommendation Code Mapping}

% For the initial specification, recommendation codes
% are assumed to consist of an age group prefix
% followed by an activity selection identifier.

% The proposed prefixes are:

% \begin{itemize}
%     \item \texttt{J}: Activities intended for ages 3--5.
%     \item \texttt{S}: Activities intended for ages 5--7.
% \end{itemize}

% For children aged five, either age group may be
% selected based on the clinician's recommendation.
% The age band boundaries and selection rules
% will be confirmed with the project supervisors.

% The following mappings are preliminary examples
% and do not represent clinically approved
% recommendations.

% \begin{longtable}{
%     p{0.22\textwidth}
%     p{0.17\textwidth}
%     p{0.49\textwidth}
% }
% \caption{Provisional Procedure Recommendation Codes}
% \label{tab:RecommendationCodes}\\

% \toprule
% \textbf{Code} &
% \textbf{Age Group} &
% \textbf{Preparation Activities} \\
% \midrule
% \endfirsthead

% \toprule
% \textbf{Code} &
% \textbf{Age Group} &
% \textbf{Preparation Activities} \\
% \midrule
% \endhead

% \texttt{J-ROUTINE} &
% 3--5 &
% Vision testing, eye drop preparation,
% and slit lamp preparation. \\

% \texttt{S-ROUTINE} &
% 5--7 &
% Vision testing, eye drop preparation,
% and slit lamp preparation. \\

% \texttt{J-VISION} &
% 3--5 &
% Vision testing preparation. \\

% \texttt{S-VISION} &
% 5--7 &
% Vision testing preparation. \\

% \texttt{J-DROPS} &
% 3--5 &
% Eye drop preparation. \\

% \texttt{S-DROPS} &
% 5--7 &
% Eye drop preparation. \\

% \texttt{J-SLIT} &
% 3--5 &
% Slit lamp examination preparation. \\

% \texttt{S-SLIT} &
% 5--7 &
% Slit lamp examination preparation. \\

% \bottomrule

% \end{longtable}

% The application will use these mappings to
% determine which activities to present when a
% valid recommendation code is entered. 
% Recommendation codes do not identify individual
% patients and are reusable across users.
% The code mappings and supported age groups
% may be revised following consultation with
% clinical stakeholders.

% \subsection{Privacy and Data Retention}

% \progname{} follows a data minimization approach, storing only
% the information necessary to support code based game
% recommendations, user preferences, and local game progress.
% The application's privacy and data management practices
% are outlined below:

% \begin{itemize}

%     \item \textbf{Data Collection:}
%     The application does not require patient accounts,
%     medical records, appointment details, or personally
%     identifying information. Optional nicknames are used
%     solely for personalization and do not need to
%     correspond to the child's legal name.

%     \item \textbf{Recommendation Codes and Clinical Data:}
%     Recommendation codes represent predefined combinations
%     of games rather than patient specific information.
%     Codes do not identify individual children or provide
%     access to clinical records. Code validation and game
%     selection are performed locally without accessing
%     clinical information systems or external services.

%     \item \textbf{Data Storage and Protection:}
%     The application stores only the active recommendation
%     code, optional user preferences, and game progress
%     required for its functionality. This information
%     is stored locally on the user's device and protected
%     from unauthorized access by other applications.
%     The application does not transmit this information
%     to external servers.

%     \item \textbf{Data Retention and Deletion:}
%     Locally stored information is retained until deleted
%     by the caregiver or through application removal,
%     subject to the device's backup and restoration
%     behaviour. Caregivers can replace or clear the
%     active recommendation code and delete stored
%     preferences and game progress.

%     \item \textbf{Consent and Access:}
%     The application does not require patient registration
%     or the collection of clinical information to access
%     its preparation games. Caregivers enter recommendation
%     codes provided by clinicians to select the games
%     relevant to their child's upcoming appointment.

%     \item \textbf{Research and Regulatory Compliance:}
%     Any personal information collected separately for
%     research or usability evaluation is outside the
%     application's operational data storage and is
%     subject to applicable institutional ethics and
%     privacy requirements.

% \end{itemize}


% ============================================================== 
\section{Non-Functional Requirements}
\label{non-func-req}

\plt{You may have to add other NFRs.}

\plt{List your nonfunctional requirements.  You may consider using a fit
  criterion to make them verifiable.}

\plt{The goal is for the nonfunctional requirements to be unambiguous, abstract
  and verifiable.  This isn't easy to show succinctly, so a good strategy may be
to give a ``high level'' view of the requirement, but allow for the details to
be covered in the Verification and Validation document.}

\plt{An absolute requirement on a quality of the system is rarely needed.  For
  instance, an accuracy of 0.0101 \% is likely fine, even if the requirement is
  for 0.01 \% accuracy.  Therefore, the emphasis will often be more on
  describing now well the quality is achieved, through experimentation, and
  possibly theory, rather than meeting some bar that was defined a priori.}

\plt{You do not need an entry for correctness in your NFRs.  The purpose of the
  SRS is to record the requirements that need to be satisfied for correctness.
  Any statement of correctness would just be redundant. Rather than discuss
  correctness, you can characterize how far away from the correct (true)
  solution you are allowed to be.  This is discussed under accuracy.}

\noindent \begin{itemize}

\item[NFR\refstepcounter{nfrnum}\thenfrnum \label{NFR_Accuracy}:]
  \textbf{Accuracy} \plt{Characterize the accuracy by giving the context/use for
    the software.  Maybe something like, ``The accuracy of the computed
    solutions should meet the level needed for $<$engineering or scientific
    application$>$.  The level of accuracy achieved by \progname{} shall be
    described following the procedure given in Section~X of the Verification and
    Validation Plan.''  If you do this provide a link to the VnV plan.}

\item[NFR\refstepcounter{nfrnum}\thenfrnum \label{NFR_Usability}:] \textbf{Usability}
  \plt{Characterize the usability by giving the context/use for the software.
    You should likely reference the user characteristics section.  The level of
    usability achieved by the software shall be described following the
    procedure given in Section~X of the Verification and Validation Plan.  Include a link to the VnV plan.}

\item[NFR\refstepcounter{nfrnum}\thenfrnum \label{NFR_Maintainability}:]
  \textbf{Maintainability} \plt{The effort required to make any of the likely
    changes listed for \progname{} should be less than FRACTION of the original
    development time.  FRACTION is then a symbolic constant that can be defined
    at the end of the report.}

\item[NFR\refstepcounter{nfrnum}\thenfrnum \label{NFR_Portability}:]
  \textbf{Portability} \plt{This NFR is easier to write than the others.  The
    systems that \progname{} should run on should be listed here.  When possible
    the specific versions of the potential operating environments should be
    given.  To make the NFR verifiable a statement could be made that the tests
    from a given section of the VnV plan can be successfully run on all of the
    possible operating environments.}

\item Other NFRs that might be discussed include verifiability,
  understandability, reusability, performance, reliability and availability, security, and accessibility.  You may have to add other NFRs.

\end{itemize}

% ============================================================== 
\section{Requirements Traceability Matrices}

\plt{The Requirements Traceability Matrix (RTM) is a table that shows the
relationship between different types of requirements. It is used to ensure that
all requirements are traceable and that no requirements are missed. The typical
flow is: Stakeholder Need -> Use Case -> Requirement}

\plt{There is also a traceability matrix between functional and nonfunctional
  requirements. The purpose of some FRs is to satisfy NFRs.  Therefore, it is
  important to show the relationship between the two.}

\plt{Further traceability to test cases and modules will be given in other
documents.}

\begin{longtable}{p{0.15\textwidth} p{0.15\textwidth} p{0.20\textwidth} p{0.20\textwidth}}

\toprule
Stakeholder Req & Use Case & Functional Req \\
\midrule
\srref{SR_Label1} & \ucref{UC_Label1} & \rref{FR_Label1} \\
\bottomrule

\caption{Traceability Between Stakeholder Requirements, Use Cases, and
Functional Requirements} \label{tab:SR2UC2FRtm}
\end{longtable}

\begin{longtable}{p{0.20\textwidth} p{0.20\textwidth}}

\toprule
Functional Req & Nonfunctional Req \\
\midrule
\rref{FR_Label1} & \nfrref{NFR_Accuracy} \\
\bottomrule
\caption{Traceability Between Functional and Nonfunctional Requirements}
\label{tab:FR2Rtm}

\end{longtable}

% ============================================================== 

\section{Rationale}

\plt{Provide a rationale for the decisions made in the documentation.  Rationale
should be provided for scope decisions, functional requirements, nonfunctioanl
requirements, data requirements, and AI requirements.  You don't need a
rationale for everything, ust the major decisions that are not obvious.  You
should have traceability to the entities in the document that you are
referencing.}

\section{Likely Changes}    

\noindent \begin{itemize}

\item[LC\refstepcounter{lcnum}\thelcnum\label{LC_meaningfulLabel}:] \plt{Give
    the likely changes, with traceability to the entity that is likely to change.}

\end{itemize}

\section{Unlikely Changes}    

\noindent \begin{itemize}

\item[LC\refstepcounter{lcnum}\thelcnum\label{LC_meaningfulLabel}:] \plt{Give
    the unlikely changes.  The design can assume that the changes listed will
    not occur.}

\end{itemize}

\newpage{}

\section{Appendices}

\subsection{Appendix A: Additional Materials}

TODO

\subsection{Appendix B: Development Plan}

\plt{This section is optional.  It is used to explain the plan for developing
  the software.  In particular, this section gives a list of the order in which
  the requirements will be implemented.  In the context of a course  this is
  where you can indicate which requirements will be implemented as part of the
  course, and which will be ``faked'' as future work.  This section can be
  organized as a prioritized list of requirements, or it could show the
  requirements that will be implemented for ``phase 1'', ``phase 2'', etc.}

\subsection{Appendix C: Values of Auxiliary Constants}

\plt{Show the values of the symbolic parameters introduced in the report.}

\plt{The definition of the requirements will likely call for SYMBOLIC\_CONSTANTS.
Their values are defined in this section for easy maintenance.}

\plt{The value of FRACTION, for the Maintainability NFR would be given here.}

TODO

\subsection{Appendix D: Generative AI Disclosure}

\textbf{AI Tool Used:} ChatGPT (OpenAI).

ChatGPT was used as a supporting tool during the preparation of
this document, primarily to assist with \LaTeX{} formatting,
improving wording and clarity, and organizing information.
Its use was mainly used in the following sections:

\begin{itemize}

    \item \textbf{Section 2.2 -- Product Vision:}
    ChatGPT was used to help summarize and organize the
    application's purpose and goals based on the team's
    existing Problem Statement and Goals document.
    The wording was refined to ensure the product vision
    remained concise, understandable, and consistent
    with the project's objectives.

    \item \textbf{Section 7 -- Data Requirements:}
    ChatGPT was used primarily to assist with formatting
    data dictionary tables in \LaTeX{}, organizing major
    data entities and their attributes, and improving
    the readability of the Privacy and Data Retention
    subsection.

    \item \textbf{General Writing and Formatting:}
    ChatGPT was used to suggest clearer wording,
    improve document organization, and identify
    potential inconsistencies or missing information.
    These suggestions were reviewed and modified
    where appropriate.

\end{itemize}

\textbf{Verification of AI-Assisted Content:}

The team reviewed AI-assisted material to ensure it was
consistent with the project's scope, existing documentation,
and intended application functionality. The team reviewed suggested 
wording, data entities, and recommendation approaches before including 
them in the document. Initial design ideas were also noted as needing 
further feedback from the project supervisors.

AI-assisted \LaTeX{} formatting was checked to ensure
that tables and other document elements were presented
correctly. The team remains responsible for the accuracy,
completeness, consistency, and appropriateness of the
final SRS document.

\subsection{Appendix E: Reflection}

Questions 1 and 2 were answered individually, while Questions 3 and 4 were
answered as a team. Questions 5 and 6 contain both team and individual
components. The team responses are provided under Team Reflection, while the
individual responses are provided under each respective team member's name.

\subsubsection*{Team Reflection}

\begin{enumerate}
    \setcounter{enumi}{2}

    \item \textbf{How many requirements were inspired by client interactions
    or their proxies?}

    A total of [number] requirements were inspired by discussions with our
    clients, stakeholders, peers, or potential users. These interactions
    primarily informed our stakeholder requirements and the functional
    requirements derived from them, helping us understand the application's
    goals, intended users, and valuable features for preparing children
    for paediatric ophthalmology appointments.

    However, we did not have sufficient time to discuss every requirement
    with stakeholders individually. We had approximately two weeks to
    complete the SRS, so we prioritized discussions about major decisions
    that would significantly affect the system if changed later.

    Consequently, many non-functional requirements were developed through
    team analysis, research, technical considerations, and assumptions
    about children aged 3--7 and their parents. In particular, our measurable
    usability and performance targets will need further validation.
    In future iterations, we would seek additional stakeholder feedback
    to verify these assumptions and refine the requirements.

    \item \textbf{Which courses have helped or will help your team
    succeed with the capstone project?}

    Several courses have provided a foundation for this project.
    SFWRENG 3RA3, which focused on software requirements, was particularly
    relevant because we previously developed a Software Requirements
    Specification and learned about requirements elicitation, documentation,
    and traceability. This experience made the structure and purpose of
    the current deliverable more familiar.

    SFWRENG 3A04, which focused on large system design, was also useful
    because it introduced concepts and activities similar to those involved
    in a capstone project. Designing a larger software system helped us
    understand the importance of defining scope, identifying system
    constraints, and considering interactions between system components.

    Additionally, SFWRENG 4HC3 was helpful because it involved interviewing
    stakeholders and eliciting requirements. These skills helped us identify
    important questions to ask our stakeholders and distinguish their needs
    from our own assumptions about the application.

    \item \textbf{What knowledge and skills will the team collectively
    need to acquire?}

    The team identified five knowledge areas essential to completing
    the capstone project:

    \begin{itemize}
        \item \textbf{Paediatric ophthalmology and clinical procedures:}
        We need to understand common paediatric eye examinations,
        including their purpose, steps, and what children may experience.
        This will help us represent procedures accurately and prepare
        children for their appointments.

        \item \textbf{Child development and child-friendly design:}
        We need to understand the cognitive, emotional, and motor
        capabilities of children aged 3--7. This will inform decisions
        about language, instructions, visual presentation, and navigation
        to ensure the interface is appropriate for the target age group.

        \item \textbf{Mobile application development:}
        We need the technical skills to implement navigation,
        interactive components, animations, audio, and local data
        storage. We must also ensure compatibility across supported
        platforms and devices.

        \item \textbf{Accessibility and usability testing:}
        We need to evaluate whether children and parents can use the
        application effectively. This includes assessing touch targets,
        text readability, colour contrast, navigation, and alternative
        ways of presenting information to identify and address usability
        barriers.

        \item \textbf{Educational game design:}
        We need to design interactive activities that reinforce learning
        objectives and familiarize children with ophthalmology procedures.
        This includes selecting appropriate game mechanics, feedback,
        and rewards that support learning while maintaining engagement.
    \end{itemize}

    These areas cover the clinical, developmental, technical, and
    educational needs of the project. Each team member will select
    at least one area based on their experience, responsibilities,
    and learning goals, and describe their plans for developing
    the relevant knowledge or skills.

    \item \textbf{What are at least two approaches to acquiring the
    knowledge or mastering each skill?}

    The team identified two approaches for each knowledge area.

    \begin{itemize}
        \item \textbf{Paediatric ophthalmology and clinical procedures:}
        \begin{enumerate}
            \item \textbf{Review clinical resources.}
            We will study reliable medical resources and educational
            materials to understand common paediatric eye examinations
            and their procedures.

            \item \textbf{Consult clinical experts.}
            We will ask clinicians and our project supervisor to clarify
            procedural details and review our understanding, helping
            us identify inaccuracies.
        \end{enumerate}

        \item \textbf{Child development and child-friendly design:}
        \begin{enumerate}
            \item \textbf{Study child development resources.}
            We will review research on the cognitive, emotional, and
            motor development of children aged 3--7 to inform decisions
            about language, instructions, visuals, and interactions.

            \item \textbf{Develop and review interface prototypes.}
            We will create screen mock-ups and navigation prototypes
            and review them with our supervisor and relevant stakeholders
            to identify confusing layouts, unclear instructions, and
            developmentally inappropriate interactions.
        \end{enumerate}

        \item \textbf{Mobile application development:}
        \begin{enumerate}
            \item \textbf{Follow official documentation and tutorials.}
            We will use the selected development framework's documentation
            and tutorials to learn how to implement navigation,
            interactive components, animations, audio, and local
            data storage.

            \item \textbf{Build small technical prototypes.}
            We will implement individual features before integrating
            them into the application to practise using the framework
            and identify technical, integration, and compatibility issues.
        \end{enumerate}

        \item \textbf{Accessibility and usability testing:}
        \begin{enumerate}
            \item \textbf{Study accessibility guidelines.}
            We will study WCAG 2.2 and relevant mobile accessibility
            guidance to evaluate text readability, colour contrast,
            touch targets, navigation, and alternative ways of presenting
            information. We will use checklists and evaluation tools
            to identify accessibility issues.

            \item \textbf{Conduct structured usability evaluations.}
            We will define tasks such as navigating to a procedure,
            completing a walkthrough, and using a practice activity.
            We will record difficulties, evaluate task completion,
            and use the findings to prioritize improvements. Where
            appropriate and approved, we will also gather feedback
            from representative users or stakeholders.
        \end{enumerate}

        \item \textbf{Educational game design:}
        \begin{enumerate}
            \item \textbf{Study educational game design principles.}
            We will review resources on game-based learning, educational
            objectives, feedback mechanisms, and age-appropriate game
            mechanics to design activities that reinforce understanding
            of ophthalmology procedures.

            \item \textbf{Prototype and evaluate gameplay.}
            We will develop small gameplay prototypes and assess whether
            their mechanics, feedback, and activity sequences support
            the intended learning objectives. Feedback from our supervisor
            and relevant stakeholders will help us identify ineffective
            mechanics and improve the connection between gameplay
            and learning.
        \end{enumerate}
    \end{itemize}

    Team members will select approaches that align with their learning
    goals and describe their plans in their individual reflections.
    \end{enumerate}

\subsubsection*{Mohammad Bilal}

\begin{enumerate}
    \item \textbf{What went well while writing this deliverable?}

    One aspect that went well was that we established a clear understanding
    of the application's purpose, scope, and features that would be excluded
    from the initial release. This provided a strong foundation for developing
    the requirements and helped the team remain aligned on the project's
    overall direction. Our project supervisor was also responsive to our
    questions, allowing us to clarify uncertainties and make decisions without
    significant delays.

    Another factor that contributed to our progress was dividing the workload
    among team members early in the process. Assigning clear responsibilities
    helped us stay organized and manage the limited time available to complete
    the deliverable. From my perspective, having a well-defined scope and clear
    responsibilities made it easier to prioritize my assigned sections, work
    efficiently, and complete them on schedule.

    \item \textbf{What pain points did you experience during this
    deliverable, and how did you resolve them?}

    One pain point was duplication and inconsistency across the
    requirements. Since multiple team members worked on different
    sections in parallel, similar ideas sometimes appeared in multiple
    places, and an idea described as an assumption in one section
    could be presented as a requirement in another. Resolving these
    inconsistencies was challenging because changes in one section
    could affect other parts of the SRS. We addressed this by keeping
    pull requests relatively small and
    having the entire team review changes that could introduce
    conflicting or duplicate information. This helped us identify
    inconsistencies, clarify the intended meaning, and remain aligned
    as the document evolved.

    Another challenge was the limited timeline. We had approximately
    two weeks to complete the SRS, compared with an entire academic
    term for a similar process in a previous requirements engineering
    course. Conducting enough stakeholder interviews to validate
    every requirement was therefore not feasible. I addressed this
    by prioritizing consultation on major decisions and making
    reasonable assumptions about smaller details when necessary.
    Although these assumptions helped us progress, they will need
    to be revisited as we receive further feedback and develop a
    stronger understanding of the project's domain.

    \setcounter{enumi}{4}

    \item \textbf{Which knowledge areas or skills will you personally
    need to acquire?}

    I will focus on accessibility and usability testing. Since the
    application is intended for children aged 3--7, it must accommodate
    their developing motor skills, limited reading ability, and varying
    accessibility needs. I need to learn how to evaluate whether the
    interface is intuitive, accessible, and easy to navigate without
    relying heavily on adult assistance. This includes assessing
    touch target sizes, text readability, colour-independent cues,
    navigation consistency, and the clarity of audio-visual instructions.
    Developing these skills will help me identify usability barriers
    and contribute to an application that is inclusive and appropriate
    for its intended users.

    \item \textbf{What approaches will you pursue to acquire this
    knowledge or master this skill, and why?}

    I will focus on studying established accessibility guidelines.
    I will review WCAG 2.2 and relevant mobile accessibility guidelines
    to understand established practices for designing accessible interfaces.
    I chose this approach because the application is intended for children
    aged 3–7, who may have limited reading ability, developing motor skills,
    and varying accessibility needs. These guidelines will help me understand
    how to evaluate important aspects of the application, such as touch
    target sizes, text readability, colour contrast, and alternative ways
    of presenting information. By learning and applying these established
    practices, I will be better equipped to identify potential accessibility
    barriers and recommend improvements to the application's design.
    This approach will also provide me with a reliable foundation for
    making informed accessibility decisions throughout development,
    rather than relying solely on personal judgment or assumptions about
    what makes an interface easy to use.

\end{enumerate}

\subsubsection*{Ahsan Muzammil}

\begin{enumerate}

    \item \textbf{What went well while writing this deliverable?}

    One thing that went well was being able to use our previous
    deliverables, such as the Problem Statement and Goals and
    Development Plan, as a foundation for the SRS. This helped
    me better understand the purpose of the application and
    how the different sections of the SRS should relate to
    one another. I also gained a better understanding of how
    requirements influence other parts of the document,
    particularly the use case scenarios, data requirements,
    and development plan. Working on these sections helped
    me think more carefully about how the application is
    expected to function and how to document these
    expectations clearly.

    \item \textbf{What pain points did you experience during this
    deliverable, and how did you resolve them?}

    One of the main challenges was maintaining consistency
    across the document and avoiding duplicated or
    contradictory information. Since different team members
    were responsible for different sections, I had to
    regularly compare my work with other sections to ensure
    that the terminology, data entities, and described
    functionality remained consistent.

    Another major challenge was the limited timeline.
    We only had approximately two weeks to develop the SRS,
    which was not enough time to confirm every requirement
    and design decision with our supervisors. Since our
    supervisors are busy medical professionals, our meetings
    mainly focused on high level decisions rather than every
    individual feature. As a result, we had to make
    assumptions and document some requirements based on
    functionality we anticipated implementing, even though
    we had not yet confirmed their feasibility or received
    approval for every detail. We addressed this by using
    our existing project goals to guide these decisions
    and identifying areas that would require further
    clarification.

    I also experienced difficulties with dependencies
    between sections. For example, I was responsible for
    Section 5 (Use Case Scenarios) and Appendix B
    (Development Plan), but I could not fully complete
    these sections until the functional and non-functional
    requirements were sufficiently developed. The use
    cases also depended on the user personas in Section 4.
    This made it difficult to work at my own pace and
    resulted in some sections being completed close to
    the deadline. I addressed this by working on the
    portions that did not depend on other sections first,
    then reviewing and updating the remaining content
    as my teammates completed their work.

    \setcounter{enumi}{4}

    \item \textbf{Which knowledge areas or skills will you personally
    need to acquire?}

    One of the main skills I need to develop is game
    development using Unity and C\#. I have never worked
    on developing a game before, and this will be my
    first time using Unity and C\# in this context.
    Although I have previous programming experience,
    I recognize that game development involves different
    concepts, tools, and workflows that I will need
    time to understand.

    In particular, I need to become familiar with
    Unity's editor, scene management, game objects,
    user interface development, animations, and
    handling user interactions. I will also need
    to improve my understanding of how C\# scripts
    are used within Unity to control game behaviour.
    Developing these skills will be important for
    contributing effectively to the implementation
    of our application's preparation games.

    \item \textbf{What approaches will you pursue to acquire this
    knowledge or master this skill, and why?}

    I plan to begin by watching beginner friendly
    Unity tutorials and following along with small
    example projects to understand the fundamentals
    of the engine. Rather than immediately working
    on complex features, I want to become comfortable
    with creating scenes, managing game objects,
    writing C\# scripts, and implementing simple
    user interactions. This will help me build a
    foundation before applying these concepts to
    our actual project.

    I also plan to use AI assisted development tools,
    particularly OpenAI Codex, to help me understand
    unfamiliar Unity concepts, troubleshoot errors,
    and learn different ways of implementing features.
    I intend to use these tools primarily as learning
    and debugging resources rather than relying on
    generated code without understanding it. By
    combining tutorials, hands on practice, AI assisted
    troubleshooting, and collaboration with my
    teammates, I hope to become more confident with
    Unity and contribute more effectively to the
    development of our application.

\end{enumerate}

\newpage

\bibliographystyle{plain}
\bibliography{../../refs/References}

\newpage


\end{document}
